\scanpage{581}

\mspage{490}

par hypothèse.

Inversement, supposons que $\underline S_{(\pi)/x}$ soit cofinal dans $M_{\pi/x}$. Comme $\deleted{\unsure{M}}$ on a

\[x\overset{\sim}{\leftarrow}\underset{M_{\pi/x}}{\unsure{\underline{\lim}|\varinjlim}}\,x_i,\]

$\pi$ étant utile pour $M$, on aura \inserted{aussi} par cofinalité

\[x\overset{\sim}{\leftarrow}\underset{\underline S_{(\pi)/x}}{\varinjlim}\,x_i,\]

or $M_{\pi/x}$ étant \unsure{grde|gde} devant $\pi$, il en est de $=$ de la sous-catégorie \deleted{\unsure{filtr.}} pleine cofinale $\underline S_{(\pi)/x}$. Donc $x$ est dans l'image essentielle de $\mathrm{Ind}_\pi\underline S_{(\pi)}$, q.e.d

\uline{Lemme} 4.16.7.12. Soient $I$ une catégorie \inserted{filtrante}, $J$ une sous-catégorie fermée, $K$ la sous-catégorie ouverte complémentaire. On a alors les équivalences

\diagram{\begin{tikzcd}[ampersand replacement=\&,column sep=large]
J\text{ cofinal}\arrow[r,Leftrightarrow]\arrow[d,Leftrightarrow] \& K\text{ non cofinal}\arrow[d,Leftrightarrow] \\
J\ne\varnothing \& K\ne I .
\end{tikzcd}}

Démonstration immédiate.

Je reviens à la démonstration de 4.16.7.8. On applique \unsure{les|le} \deleted{deux} \unsure{lemmes|lemme} 4.16.7.11 \deleted{\unsure{aux deux|aux de}} à $M,\pi$ et aux deux

% archive sheet number: 580
